\clearpage
\fancyhead[L]{\small BRAMBILLA. EFFEMERIDI ASTRONOMICHE.}
\SourceLeaf{634}
\begin{center}\small Göttingische gelehrte Anzeigen. No. 64. pp. 633..637. 21 April 1823.\\[3pt]\rule{.72\linewidth}{.3pt}\end{center}

\emph{Effemeridi Astronomiche di Milano per l'anno \textup{1823} calcolate da} \textsc{Enrico Brambilla}. \emph{Con Appendice.} Milan, 1822. The ephemerides themselves comprise 112 pages, and the appendix 88 pages, in octavo.

\begingroup\baselineskip=12.5pt\relax
It has long been customary for the Milan ephemerides to include several essays whose value extends beyond the year of publication and which make the ephemerides interesting even to people who have no occasion to use the astronomical calendar.
All we need say about the calendar is that its commendable arrangement has remained unchanged. The appendix contains the following essays.

Geographical position of some mountains visible from Milan, by \textsc{Barnabas Oriani}.
Here this distinguished astronomer gives us a list of 49 points in northern Italy and Switzerland, with their latitudes and longitudes, their positions relative to Milan Cathedral, and their elevations above sea level.
These determinations are based for the most part on measurements that the Milan astronomers carried out in 1788---1791 by order of the Austrian government and in 1803---1806 by order of the Italian government.
The measurements did not extend as far as the sea, so the elevations were only relative. The barometric determination of Milan's elevation served as the basis for reducing them to sea level.
Some determinations in the list are based on recent surveys, which are described here in somewhat greater detail and are of particular interest.
The position of Monte Viso in Piedmont was determined by measurements taken from Turin and Milan. It lies 96380 toises (25 geographical miles) from Milan.
The elevation measurements taken from the two places agree best when terrestrial refraction is taken as 0.08 times the curvature of the terrestrial arc, and \textsc{Oriani} accordingly puts the elevation above sea level at 1968 toises.
The determination of the elevation of Monte Cimone in the Apennines near Lucca at 1112 toises is particularly noteworthy because it is based on direct measurement of the dip of the sea horizon for both the Mediterranean and the Adriatic. \textsc{Brioschi} carried out these measurements in 1817.
The elevation of Monte Rosa is put at 2385, that of Finsterarhorn at 2203, and that of the Simplon at 1805 toises. --- Report on the operations carried out in 1822 to determine the differences in longitude between several places in Italy by means of gunpowder signals on Monte Cimone, by \textsc{Francesco Carlini}.
Once the triangulations in France and the Austrian states have been completed and linked, they will furnish the measurement of a very large arc of the middle parallel, extending for about 24 degrees from the Atlantic Ocean to Orsova.
It is obvious how important these measurements, when combined with suitable astronomical operations, may become for a more complete knowledge of the figure of the Earth.
The terrain itself through which this arc runs is particularly favorable for the operations by which the difference in longitude between the endpoints must be determined. Along this stretch there are so many high mountains commanding an immensely wide view that a comparatively very small number of intermediate points will suffice.
An initial trial of this kind was already made in Sept. 1821, when the difference in longitude between the Milan observatory and the hospice on Mont Cenis was determined by means of gunpowder signals on Rocca Melone, which is 1792 toises high and 86000 toises from Milan.
Encouraged by the successful outcome of this trial, they wished to proceed to a
\SourceLeaf{635} more extensive connection.
To that end the French geographers put forward a bold plan. They thought that by using three intermediate points, namely the aforementioned Monte Viso, Monte Cero near Padua, and Monte Maggiore in Friuli, they could establish the connection between the east coast of the Adriatic Sea and Mont d'Or near Clermont in the middle of France in a single night.
They backed away from it, however, because they considered the difficulties too great.
It is still uncertain whether Monte Viso can be climbed at all, and still more uncertain whether one can stay on its highest, steep summit during the night.
And even supposing these difficulties could be overcome, they fear that at the enormous distance from Monte Cero (50 geogr. miles) the light of the gunpowder flashes would remain invisible even through the most powerful telescopes
(These remarks suggest that this immense plan has not yet been abandoned outright. From the stated heights of these two points and their distance apart, however, the reviewer finds that neither can possibly rise above the other's physical horizon. Were it not for this circumstance, the last two difficulties could be addressed by using large heliotropes).
They therefore decided for the time being on a more limited operation, having gunpowder signals given on Monte Cimone over several nights at the beginning of May 1822, to be observed from Monte Cero, in Milan, and at various other Italian observatories.
But unfavorable weather prevented these operations from succeeding in their main objective. The gunpowder signals were seen neither in Milan nor on Monte Cero.
They were observed, however, in Parma, Modena, Bologna and Florence. Yet the differences, in some cases considerable, between the results and the differences in longitude yielded by the geodetic operations seem to prove that the necessary precision in determining the time was lacking at some of these places.
What is particularly noteworthy, moreover, is that Mr.~\textsc{Carlini}, from his observations with an 18-inch \textsc{Reichenbach} repeating circle, found the altitude of the pole at Parma to be $22''6$ greater than that yielded by the geodetic connection with Milan.
Since similar anomalies had already appeared at several other places in northern Italy, it is difficult to doubt their reality, and it is much to be wished that all the principal triangles of the Austrian triangulations should soon be made known in full. --- The remaining articles in the appendix also contain the oppositions of Uranus observed by \textsc{Angelo Cesaris} in 1817 and 1818 and the oppositions of Jupiter and Saturn observed by the same astronomer in 1821. They also contain occultations of stars and eclipses of Jupiter's satellites observed by \textsc{Hallaschka} in Prague and his observations of the first comet of 1822. Finally, they contain the meteorological observations in Milan in 1820 by A. \textsc{Cesaris}.
\par\clearpage\endgroup
